\documentclass[%
 reprint,
%superscriptaddress,
%groupedaddress,
%unsortedaddress,
%runinaddress,
%frontmatterverbose, 
%preprint,
%preprintnumbers,
%nofootinbib,
%nobibnotes,
%bibnotes,
 amsmath,amssymb,
 aps,
%pra,
%prb,
%rmp,
%prstab,
%prstper,
%floatfix,
]{revtex4-2}


\usepackage[utf8]{inputenc} % allow utf-8 input
\usepackage[T1]{fontenc}    % use 8-bit T1 fonts
\usepackage{hyperref}       % hyperlinks
\usepackage{url}            % simple URL typesetting
\usepackage{booktabs}       % professional-quality tables
\usepackage{amsfonts,amsmath,amssymb}       % blackboard math symbols
\usepackage{nicefrac}       % compact symbols for 1/2, etc.
\usepackage{microtype}      % microtypography
\hypersetup{colorlinks=true,linkcolor=blue,citecolor=red!60!blue}
\usepackage[cmyk]{xcolor}
\usepackage{graphicx}

\usepackage{caption}
\captionsetup{justification=raggedright,singlelinecheck=true}
\usepackage{subcaption}

\newtheorem{thm}{Theorem}
\newtheorem{cor}{Corollary}
\newtheorem{prop}{Proposition}
\newtheorem{lem}{Lemma}

% \theoremstyle{definition}
\newtheorem{definition}{Definition}
\newtheorem{method}{Method}

\DeclareMathOperator{\tr}{Tr}
\DeclareMathOperator{\Tr}{Tr}

\newcommand{\argmax}[1]{\underset{#1}{\operatorname{arg}\,\operatorname{max}}\;}
\newcommand{\argmin}[1]{\underset{#1}{\operatorname{arg}\,\operatorname{min}}\;}

\newcommand{\norm}[1]{\lVert{#1}\rVert}

\renewcommand{\a}{{\bf a}}
\renewcommand{\b}{{\bf b}}
\renewcommand{\c}{{\bf c}}
\renewcommand{\d}{{\bf d}}
\newcommand{\g}{{\bf g}}
\newcommand{\h}{{\bf h}}
\renewcommand{\k}{{\bf k}}
\newcommand{\m}{{\bf m}}
\newcommand{\n}{{\bf n}}
\newcommand{\p}{{\bf p}}
\newcommand{\s}{{\bf s}}
\newcommand{\w}{{\bf w}}
\newcommand{\x}{{\bf x}}
\newcommand{\y}{{\bf y}}
\newcommand{\z}{{\bf z}}

\newcommand{\A}{{\bf A}}
\newcommand{\B}{{\bf B}}
\newcommand{\C}{{\bf C}}
\newcommand{\D}{{\bf D}}
\newcommand{\E}{\mathbb{E}}
\newcommand{\F}{{\bf F}}
\newcommand{\G}{{\bf G}}
\renewcommand{\H}{{\bf H}}
\newcommand{\I}{{\bf I}}
\newcommand{\K}{{\bf K}}
\renewcommand{\L}{{\bf L}}
\newcommand{\M}{{\bf M}}
\newcommand{\N}{{\bf N}}
\renewcommand{\P}{{\bf P}}
\newcommand{\Q}{{\bf Q}}
\newcommand{\R}{\mathbb{R}}
\renewcommand{\S}{{\bf S}}
\newcommand{\U}{{\bf U}}
\newcommand{\V}{{\bf V}}
\newcommand{\W}{{\bf W}}
\newcommand{\X}{{\bf X}}
\newcommand{\Y}{{\bf Y}}
\newcommand{\Z}{{\bf Z}}

\newcommand{\Gam}{\boldsymbol{\Gamma}}
\newcommand{\Thet}{\boldsymbol{\Xi}}

\newcommand{\thet}{\boldsymbol{\xi}}

\renewcommand{\r}{{\bf c}^a}
\renewcommand{\v}{{\bf v}}
\renewcommand{\u}{{\bf u}}
\newcommand{\e}{{\epsilon}}

\renewcommand{\H}{{\bf H}}
\renewcommand{\P}{{\bf P}}
\renewcommand{\S}{{\bf S}}

\newcommand{\Real}{{\mathbb R}}

\newcommand{\Lam}{\boldsymbol{\Lambda}}
\newcommand{\Sig}{\boldsymbol{\Sigma}}

\newcommand{\lb}{\left(}
\newcommand{\rb}{\right)}

\newcommand{\wt}{\widetilde}
% \renewcommand{\ol}{\overline}
\newcommand\note[1]{\textcolor{red}{\uppercase{[#1]}}}

\begin{document}

\title{Neuronal Temporal Filters as Normal Mode Extractors}




% \author{Siavash Golkar\thanks{Equal contribution}\hspace{3pt}$^{\,1}$    \hspace{25pt}   Jules Berman\protect\footnotemark[1]\hspace{3pt}$^{\,1}$  \hspace{25pt}   Shiva Farashahi$^{\,1}$ \vspace{7pt}
%   \\
%   \textbf{David Lipshutz$^{\,1}$ \hspace{25pt} Dmitri B.\ Chklovskii$^{\,1,2}$  }
%   \vspace{14pt}
%   \\
%   $^{1\,}$Center for Computational Neuroscience, Flatiron Institute
%   \\
%   $^{2\,}$Neuroscience Institute, NYU Medical Center
%   \vspace{5pt}\\
%   \texttt{\{sgolkar,jberman,sfarashahi,dlipshutz,dchklovskii\}@flatironinstitute.org}
% }



\author{Siavash Golkar}
\thanks{Equal contribution}
\affiliation{Center for Computational Neuroscience, Flatiron Institute}
\author{Jules Berman}
\thanks{Equal contribution}
\affiliation{Center for Computational Neuroscience, Flatiron Institute}
\author{David Lipshutz}
\affiliation{Center for Computational Neuroscience, Flatiron Institute}
\author{Robert Mihai Haret}
\affiliation{University Medical Center Göttingen, Department of Ophthalmology and Bernstein Center for Computational Neuroscience Göttingen}
\author{Tim Gollisch}
\affiliation{University Medical Center Göttingen, Department of Ophthalmology and Bernstein Center for Computational Neuroscience Göttingen}
\author{Dmitri B.\ Chklovskii}
\affiliation{Center for Computational Neuroscience, Flatiron Institute}
\affiliation{Neuroscience Institute, NYU Medical Center}
   
\date{\today}% It is always \today, today,
             %  but any date may be explicitly specified



\begin{abstract}
To generate actions in the face of physiological delays, the brain must predict the future. Here we explore how prediction may lie at the core of brain function by considering a neuron predicting the future of a scalar time series input. Assuming that the dynamics of the lag vector (a vector composed of several consecutive elements of the time series) are locally linear, Normal Mode Decomposition  decomposes the dynamics into independently evolving (eigen-)modes allowing for straightforward prediction. We propose that a neuron learns the top mode and projects its input onto the associated subspace. Under this interpretation, the temporal filter of a neuron corresponds to the left eigenvector of a generalized eigenvalue problem. We mathematically analyze the operation of such an algorithm on noisy observations of synthetic data generated by a linear system. Interestingly, the shape of the temporal filter varies with the signal-to-noise ratio (SNR): a noisy input yields a monophasic filter and a growing SNR leads to multiphasic filters with progressively greater number of phases. Such variation in the temporal filter with input SNR resembles that observed experimentally in biological neurons.
% Next, we apply our algorithm to natural stimuli that are generated by nonlinear dynamics and obtain temporal filters that resemble those of biological neurons. Remarkably, these filters derived in our unsupervised algorithm, predict the neural output on par with the supervised linear kernel trained to perform this task.
\end{abstract}

\maketitle

%\tableofcontents
% Outline:
% 1.	Introduction
% a.	Modeling neuron operation with a prediction algorithm
% b.	Prior work: predictive coding and predictive information
% c.	Linear dynamics prediction – NMD
% d.	Contributions/relationship with prior work
% e.	Paper organization
% 2.  Problem statement
% 3.	Detecting dominant mode in a noisy sum of exponentials
% a.	Problem formulation/lag vector dynamics (Figure 1: three panels from exponentials to the sum of exponentials to the dominant mode)
% b.	DMD solution: eigendecomposition
% c.	Dependence of the number of phases on the rank (Figure 2: filter as a function of exponents, rank/SNR, recovery of the mode – eigenvalue and exponential)
% d.	Impact of regularization (Figure 3)
% 4.	Learning on concatenated synthetic functions 
% a.	Concatenated sums of exponential (Figure 4a: Concatenated sums of exponential, filter shape dependence on rank/SNR)
% b.	Logistic functions (Figure 4b: Concatenated logistic function, filter shape, dependence on the time constant and rank/SNR)
% 5.	Learning on actual data
% a.	LMC data results (Figure 5a: Data and filter shapes)
% b.	Comparison with supervised learning (Figure 5b: comparison with supervised learning filter)
% 6.	Discussion
% a.	Relevance of linear NMD to nonlinear dynamics
% b.	New explanation of the temporal filter
% 7.	Supplementary data
% a.	Controls

 
\section{Introduction}

The brain must generate behavior in the face of physiological delays on multiple levels: from sensory transduction to axonal conduction to synaptic transmission and muscle activation. To compensate for such delays, it would be useful even for individual neurons to predict future inputs. Not surprisingly, the paradigm of optimal prediction has been used to derive normative models of neurons. Such models can be loosely classified into two categories \citep{singer2018sensory,chalk2018toward}: predictive coding \citep{srinivasan1982predictive,rao1999predictive,huang2011predictive, gregor2012lattice} and predictive information \citep{tishby2000information,bialek2006efficient}. 

In predictive coding, a neuron computes an optimal prediction of the signal and subtracts it from the actual value, transmitting prediction error downstream \citep{srinivasan1982predictive,rao1999predictive,huang2011predictive}. This results in an optimal linear filter whose shape depends on the statistics of the input. Predictive coding approximately accounts for the change in the temporal receptive field with the input SNR (Fig.~\ref{fig:intro}A) but fails to reproduce the exact shape of the filter because it has a narrow unitary peak at the exact time of the actual signal (usually present time) \citep{srinivasan1982predictive}. Such peak can be smoothed by adding a de-noising objective but that introduces additional parameters to the model. 

In predictive information, a neuron filters the past signal to preserve only the part which is most informative about the future \citep{tishby2000information,bialek2006efficient,chalk2018toward}. Such framework expands the assortment of temporal filters depending on the various parameters such as the balance between the past/output and output/future mutual information. Recently, it was shown that multi-phasic filter can be derived using this formalism as well~\cite{retinal-mosaics}.

A shortcoming of both approaches is that they require setting the time point relative to the present for which the neuron generates prediction. Such formulation of the prediction problem does not seem suitable when neurons might have different amounts of delay, or the goal is to predict (and act upon) a trend or a mode of the signal that spans multiple time points in the future. 
Moreover, both approaches view the neuronal filter as a trade off between different mathematical terms in the optimization objective and therefore depend on the relative weighting of these terms. Therefore, the shape of the temporal kernel depends on these - as well as other - hyperparameters of the model.
% Finally, neither approach seems to accurately reproduce the multi-phasic nature of the temporal kernel in different noise conditions \citep{mainen1995reliability} (Fig.~\ref{fig:intro}A).

%An alternative approach to prediction has been developed under the name of Dynamic Mode Decomposition (DMD) \citep{schmid2010dynamic,kutz2017dynamic}. This approach assumes that a generative model for the input signal is a linear dynamical system, although it has also been successfully applied to complex nonlinear systems \citep{schmid2010dynamic,kutz2017dynamic}. The NMD algorithm decomposes the input into the eigenmodes that evolve independently of each other allowing a straightforward parameter-free prediction of the future at various time points.
% 
\begin{figure}[h]
     \centering
    \includegraphics[width=0.42\textwidth]{finalfigures/mono-bi.pdf}
    \includegraphics[width=0.42\textwidth]{finalfigures/bi-tri.pdf}
    \includegraphics[width=0.42\textwidth]{finalfigures/vsnoise.pdf}
    \label{fig:samples}
    % \vspace{-2pt}
        \caption{Neuronal temporal filters. (A,B) Experimentally obtained by Spike Triggered Average (STA) from two retinal ganglion cells stimulated by white noise of different contrast. (C) Top left eigenvectors obtained in our framework for different levels of noise in a synthetic example. Shading shows standard deviation.  Whereas theoretical filters peak when the neuron spikes, experimental ones peak before that. This artifact is due to signal processing delays from photo-absorption to the ganglion cell firing and the smoothness of experimental filters resulting from the low-pass nature of biophysical processes.}
    \label{fig:intro}
\end{figure}
An alternative approach to prediction assumes that the sensory stimuli are generated by a (generally nonlinear) dynamical system. If such dynamics can be learned by the brain from previously seen trajectories, then prediction can be cast as identifying long-living (growing or slow) modes (or manifolds) in the input and extrapolating them into the future. To take the first step in this direction, we take advantage of the fact that in the vicinity of a hyperbolic fixed point the invariant manifolds can be well approximated by the invariant subspaces of the linearized dynamics (see e.g. Sec.~19.12a of~\cite{wiggins2006introduction}). Therefore, a neuron that identifies the least decaying linear mode can identify an invariant manifold of a fixed point and generate a non-trivial prediction.

Here, we explore the hypothesis that a neuron extracts the least decaying mode from the incoming scalar time series as in Normal Mode Decomposition \citep{Landau1976Mechanics}. To cast this problem in the language of dynamics, we construct time-lag vectors by windowing the scalar time series. We assume that the system is in the vicinity of a hyperbolic fixed point and identify the linear dynamics of the lag vectors. The eigen-decomposition of the linear transition matrix yields the eigenmodes that evolve independently as the right eigenvectors and the corresponding linear filters as the left eigenvectors. Therefore, the top left eigenvector yields a linear filter that projects onto the most dominant mode, representing neuronal output. Thus, unlike the previous predictive coding and predictive information approaches which optimize prediction for a certain time in the future, in our proposal, a neuron learns a rank-1 approximation of the dynamical system generating the input and outputs the top mode which develops into the future in a predictable way.

Our theoretical proposal makes a prediction that can be compared with experimental data. We find that the shape of the top left eigenvector depends on the input SNR: monophasic filter at low SNR with adding phases for growing SNR (Fig.~\ref{fig:intro}C). This prediction is reminiscent of the SNR-induced variation in the form of the Spike-Triggered Average (STA) of various types of biological neurons (Fig.~\ref{fig:intro}AB)\cite{srinivasan1982predictive,baccus2002fast,liu2015spike,chander2001adaptation,kim2001temporal,mainen1995reliability}. Although this does not prove our proposal it suggests that this approach may be a step in the right direction.

%Here, we explore a model of a neuron receiving a scalar time series, performing NMD on the input and outputting the dominant mode. We assume that the time series is generated by a dynamical system with a hyperbolic fixed point, and that the system is in the vicinity of this fixed point. While this is a simplifying assumption, it provides a useful processing step, as non-linear systems often comprise multiple fixed point which then act as a skeleton for different types of trajectories that traverse the phase space between them~\cite{?}. In a similar way, neurons that specialize to learning the dynamics near the different fixed points of the non-linear system can then act as a foundation for further processing downstream. In this work we lay the groundwork for this approach by considering the NMD near one of these fixed points. This implies that  the system is dominated by the properties of the linearized dynamics and in particular performing NMD is equivalent to projecting on the eigenvectors given by the largest eigenvalues of the hyperbolic fixed point. 
% 
%Mathematically, such model reduces to solving a generalized eigenproblem defined on the incoming data expressed as Hankel matrices. The temporal filter of a neuron is given by the top left eigenvector which can be computed straightforwardly. Interestingly, the shape of the top left eigenvector depends on the input SNR: monophasic filter at low SNR with adding phases for growing SNR (Fig.~\ref{fig:intro}B). This result accounts for experimental observations on biological neurons. 

%Unlike the previous predictive coding and predictive information approaches which optimize prediction for a certain time in the future, the NMD neuron learns a rank-1 approximation of the dynamical system generating the input and outputs the top mode which develops into the future in a predictable way. Moreover, NMD yields parameter-free filters. 

% Finally, we apply our framework to electrophysiological experiments on visual neurons. Our unsupervised approach predicts measured output from the input at least as well as the optimal linear filter trained in a supervised setting does (on the test data). This suggests that our theory captures an essential part of the neuronal function.

%A central problem in neuroscience is to understand, on the algorithmic level, the transformation performed by a neuron on its inputs. Whereas biophysical models can fit diverse experimental data they do not illuminate the computational role of a neuron. Simpler models, such as rectifying linear (or linear-nonlinear) units for graded potential neurons and integrate-and-fire units for spiking neurons are commonly used, but also do not provide sufficient insight. 

%Neuronal models are usually learned in the supervised setting by minimizing the mismatch between predicted and observed outputs for a given input. This approach is natural in fitting a neuronal model but again does not offer insight into its function. Indeed, typical neurons lack a supervision signal and operate in a regime closer to unsupervised setting. Therefore, truly understanding the functional role of a neuron may require modeling it as an unsupervised learning algorithm.

%To test our predictions, we compute optimal filters for the experimentally generated neuronal input for different levels of noise. We show that the filter varies with the level of noise as observed experimentally. We predict the neuronal output by projecting its input on the optimal filter and demonstrate that the prediction error in our unsupervised formalism is comparable to competing supervised algorithms. 
% \clearpage
\section{Normal Mode Decomposition (NMD)}\label{sec:DMD}
In this Section, we review NMD \citep{Landau1976Mechanics} in the context of a neuron with a Single Input and a Single Output (SISO)~\footnote{In a neuron receiving synapses from multiple sources this would correspond to considering the total synaptic current as a scalar input.}. 

\subsection{Problem statement and linearization}

Let $\{x(t)\}_{t=1,2,\dots}$ be a scalar time series satisfying the non-linear relation
\begin{equation}\label{eq:nl_dyn}
    x(t+1)=f(x(t),\dots,x(t-n)),\qquad t\ge n+1,
\end{equation}
where $f:\R^n\to\R$ is a continuously differentiable function and $n\ge1$ is referred to as the \textit{lag}. We say $x^\ast\in\R$ is a \textit{fixed point} if $x^\ast=f(x^\ast,\dots,x^\ast)$. We embed a scalar time series $x(t)$, $t=1,\dots,T+n$, into a sequence of lag vectors $\x_t=[x(t-n+1),\dots,x(t)]^\top\in\R^n$. The dynamics of the first $n-1$ components of the lag vector is given by a simple shift in the time-step (i.e. $x(t-i)\to x(t-i+1)$) and the dynamics of the final component $x(t)$ is given by Eq.~\eqref{eq:nl_dyn} (i.e. $x(t)\to f(x(t),\dots,x(t-n))$). These can be summarized as $\x_{t+1}=\F(\x_t)$, where the function $\F$ is given by
\begin{equation}
\label{dyn_nl}
    \F(\x_t) = \F(\begin{bmatrix}
        x(t-n+1) \\ \vdots \\ x(t-1) \\ x(t)
    \end{bmatrix})\equiv\begin{bmatrix}
        x(t-n+2) \\ \vdots \\ x(t) \\ \\ f(x(t),\dots,x(t-n))
    \end{bmatrix}
\end{equation}
We assume that the dynamics are dominated by the presence of a hyperbolic fixed point at $\x^*$. 
We approximate dynamics \eqref{dyn_nl} by expanding around this fixed point. Defining $\delta \x_t = \x_t - \x^*$, we have:
\begin{align*}
    \x_{t+1}&=\F(\x^*)+ \nabla\F(\x^*)\delta\x_t+o(\|\delta\x_t\|)\\
    &\approx \x^* + \nabla\F(\x^*)\delta\x_t,
    % =\nabla\F(\x^*)\x_t+\F(\x^*)-\nabla\F(\x^*)\x^*,
\end{align*}
where we have used the fact that $\F(\x^*)=\x^*$ and $\nabla\F(\x^*)$ is the Jacobian evaluated at $\x^*$. After subtracting $\x^*$ from both sides:
\begin{align*}
    \delta \x_{t+1}\approx\nabla\F(\x^*)\delta\x_t
\end{align*}
Letting $\A:=\nabla\F(\x^*)$ and assuming without loss of generality that $\x^*=
\bf0$, we have the following linear approximation to the nonlinear dynamics in equation \eqref{dyn_nl} near the fixed point:
\begin{equation}
\label{dyn}
    \x_{t+1}=\A\x_t.
\end{equation}
From the definition of the function $\F$ in Eq.~\eqref{dyn_nl}, we see that $\A=\nabla\F(\bf 0)$ takes a companion matrix form~\citep{bellman1997introduction}
\begin{equation}\label{eq:A_def}
    \A=
    \begin{bmatrix}
        0 & 1 &  &   \\
         & 0 & 1 &   \\
         &  & \ddots & \ddots &   \\
         &  &  &  0 & 1 \\
        c_1 & c_2 & \hdots & &c_n
    \end{bmatrix}.
\end{equation}
where $c_i=\partial f(\x^\ast)/\partial x^i$.

\subsection{Eigendecomposition of the dynamics}
The matrix $\A\in \R^{n\times n}$ is generically a non-normal matrix possessing an eigendecomposition:
\begin{equation}
\label{DMD}
    \A=\sum_i \lambda_i\w_i \v_i^\top \;\;\; {\rm s.t.} \;\;\; \v_i^\top\w_j=\delta_{ij},
\end{equation}
where $\lambda_i$ are the eigenvalues and $\w_i$ (resp. $\v_i^\top$) are the right (resp. left) eigenvectors.
% Whereas, typically, NMD is applied to oscillatory dynamics with complex $\lambda_i$, 
Here, we consider
 matrix $\mathbf{A}$ describing a hyperbolic fixed point ($\lambda_i\neq 1$). Moreover, we assume, that all eigenvalues are real and are sorted in descending order \mbox{$\lambda_i\geq \lambda_{i+1}$}. We also assume that the top mode is non-degenerate (i.e. $\lambda_1>\lambda_2$).  


% If $\b\neq 0$, Eq.\eqref{mse} can be optimized by taking advantage of the fact that $\x_t$ and $\x_{t+1}$ in Eq.\eqref{dyn} are lag vectors and 

% \begin{align*}\label{eq:companion}
%     \A=
%     \begin{bmatrix}
%         0 & 1 &  &   \\
%          & 0 & 1 &   \\
%          &  & \ddots & \ddots &   \\
%          &  &  &  0 & 1 \\
%         c_1 & c_2 & \hdots & &c_n
%     \end{bmatrix}.
% \end{align*}
% Then the error must be optimized only with respect to vectors $\b$ and $\c$.
% \begin{equation}
% \label{msec}
%     \min_{\c,\b}\sum_t\left\lVert\x_{t+1}-\A\x_t-\b\right\rVert^2.
% \end{equation}
% Then, we perform the eigendecomposition of $\A$. 
The normal modes are found by eigendecomposition of~$\A$.  As a matrix of companion form, $\A$ is diagonalized using the Vandermonde matrix and its inverse~\citep{bellman1997introduction}:
\begin{align*}
    \A=\V\Lam\V^{-1},
\end{align*}
where $\V$ is the Vandermonde matrix
\begin{align*}
    \V:=\begin{bmatrix}
        1&1&\cdots&1\\
        \lambda_1&\lambda_2&\cdots & \lambda_n\\
        \vdots& \vdots & & \vdots\\
        \lambda_1^{n-1} & \lambda_2^{n-1} & \cdots & \lambda_n^{n-1}
    \end{bmatrix}
\end{align*}
and $\Lam:=\text{diag}(\lambda_1,\dots,\lambda_n)$ is the diagonal matrix of eigenvalues $\lambda_i$, which are the roots of the characteristic polynomial with coefficients $c_i$. The left eigenvectors of $\A$, are given by the rows of the inverse of the Vandermonde matrix $\V^{-1}$. It is easily verified that the right eigenvectors of $\A$, corresponding to the columns of the Vandermonde matrix, represent the individual modes of the dynamical system.  

\subsection{Projecting onto the dominant mode}
We are interested in finding the dominant exponential mode of the dynamics. As stated above, the dominant exponential mode is represented as the top right eigenvector of $\A$. Therefore, by projecting onto the top right eigenvector, we can isolate the dominant mode. 

Because of the bi-orthogonality of the left and the right eigenvectors, the top left eigenvector (henceforth referred to as $\v_1$) is orthogonal to all but the top right eigenvector. This allows us to use $\v_1$ as a projector that zeros out all but the most dominant mode of the dynamics. We therefore propose that it is the neuron's goal to learn the top left eigenvector and project its input onto this vector.

In order to learn the top left eigenvector, one approach would be to first find matrix $\A$ and then perform the eigen-decomposition on the inferred $\A$ matrix. This matrix can be found from data $x(t)$ by minimizing the mean squared error:
\begin{equation}
\label{mse}
    \min_{\A}\sum_t\left\lVert\x_{t+1}-\A\x_t\right\rVert^2.
\end{equation}
Eq.~\eqref{mse} can be solved via the following system of equations:
\begin{equation}
\label{mse-sol}
  \X_+\X^\top = \A\X\X^\top,
\end{equation}
where we introduce a matrix notation $\X=[\x_n,...,\x_{T+n-1}]$ and $\X_+=[\x_{n+1},...,\x_{T+n}]$.

In this paper, however, instead of directly solving for $\A$ via Eqs.~\eqref{mse} and \eqref{mse-sol}, we substitute the eigendecomposition, Eq.\eqref{DMD} and multiply both sides by $\v_i^\top$ on the left, leading to the equivalent generalized eigenvalue problem:
\begin{equation}\label{eq:gen_eig_prob}
  \v_i^\top \X_+\X^\top = \lambda_i\v_i^\top\X\X^\top.
\end{equation}
This allows us to circumvent the intermediate step of computing the $\A$ matrix explicitly.


% The logic is that the decaying exponentials and oscillatory signals have been already seen by the brain in the recent past. Thus, they could have been acted upon and do not need continued attention.  Growing exponentials, on the other hand, represent nascent potentially actionable trends. 


% In the SISO case, we can extract only one exponential. Because the greatest contribution to the future is given by the fastest growing  mode, we compute the top left eigenvector, $\v_1$, corresponding to the top eigenvalue, $\lambda_1$, as judged by its real part, and project the input onto $\v_1$. Yet the form of $\v_1$ depends on the other right eigenvectors of $\A$ because of their bi-orthogonality, see equation \eqref{DMD}. As the right eigenvectors corresponding to real eigenvalues are exponential, bi-orthogonality implies that the top left eigenvector changes sign $n-1$ times for $k=n$ (see SM Sec.~\ref{app:vandermonde}). %Therefore, the left eigenvectors of $\A$ have $k$ phases.

%In a data-driven algorithm, ideally we would compute the elements $c_i$ of the companion matrix $\A$ and use this to compute the left eigenvectors. However, imposing this structure on the form of $\A$ can be challenging in a biological algorithm. Therefore, in what follows, we will compute $\A$ from data without any constraints on the form of $\A$ with the possible side-effect of increasing sensitivity to noise. In practice, instead of solving for $\A$ from data, we solve the generalized eigenvalue problem:
%\begin{equation}
 %   \v_i\A=\lambda_i \v_i \Rightarrow \v_i\langle\x_{t+1}\x_t^\top\rangle = \lambda_i\v_i \langle \x_t \x_t^\top\rangle
%\end{equation}
%This avoids deducing $\A$ via $\langle\x_{t+1}\x_t^\top\rangle\langle \x_t \x_t^\top\rangle^{-1}$.


% \section{Connection to Singular Value Decomposition}

% In the previous section we argued that the problem of dominant mode detection can be reduced to that of finding the top left eigenvector $v_1$. Here, we will show that this problem is very closely related to a Singular Value Decomposition (SVD) problem. This relationship is of interest, in part because finding the top singular vector (epsecially when applied to non-normal matrices) is a more numerically stable problem compared to finding the top eigenvector. Therefore, by establishing this connection we show that the original eigenvector problem is both well-defined and easy to find numerically.

% To see the relationship between our eigenvalue formulation and SVD, we first look at the eigendecomposition of $\A^m$ where $m$ is an integer number:
% \begin{equation}
%     \A^m = \left(\sum_{i=1}^n \lambda_i \u_i \v_i^\top \right)^m=\sum_{i=1}^n \lambda_i^m \u_i \v_i^\top.
% \end{equation}
% If we assume that the dominant mode is real and is non-degenerate, as we take $m\to\infty$, $\A$ will become approximately a rank one matrix:
% \begin{equation}
%     \lim_{m\to\infty}\A^m = \lambda_1^m \u_1 \v_1^\top + \mathcal O\left(\Big(\frac{\lambda_2}{\lambda_1}\Big)^m\right).
% \end{equation}
% Therefore, in this limit, the SVD of the matrix $\A^m$ will coincide with the eigen-decomposition. In particular, the right singular vector of $\A$ will be $\v_1^\top$, the same as the left eigenvector. We show this empirically in (Fig.~\ref{fig:svd_cong}) where we see the convergence of of the right singular vector of $\A^m$ to the left eigenvector of $\A$ for increasing $m$.

% \begin{figure}[h]
%      \centering
%      \begin{subfigure}[b]{0.23\textwidth}
%          \centering
%          \includegraphics[width=\textwidth]{svd/RSV_TO_LEV.png}
%      \end{subfigure}
%      \hfill
%      \begin{subfigure}[b]{0.23\textwidth}
%          \centering
%          \includegraphics[width=\textwidth]{svd/L2.png}
%      \end{subfigure}
%         \caption{Convergence of Right Singular Vector to Left Eigenvector for powers of A. Is is computed on time series made of random sums of exponentials.}
%         \label{fig:svd_cong}
% \end{figure}


% \begin{figure}[h]
%      \centering
%      \begin{subfigure}[b]{0.23\textwidth}
%          \centering
%          \includegraphics[width=\textwidth]{svd/LEV_NOISE.png}
%      \end{subfigure}
%      \hfill
%      \begin{subfigure}[b]{0.23\textwidth}
%          \centering
%          \includegraphics[width=\textwidth]{svd/RSV_NOISE.png}
%      \end{subfigure}
%         \caption{Noise response of Left Eigenvector and Right Singular Vector}
%         \label{fig:svd_noise}
% \end{figure}


% \begin{figure*}[ht]
%      \centering
%      \begin{subfigure}[b]{0.49\textwidth}
%          \centering
%          \includegraphics[width=\textwidth]{svd/CONJ_FALSE.png}
%      \end{subfigure}
%      \hfill
%      \begin{subfigure}[b]{0.49\textwidth}
%          \centering
%          \includegraphics[width=\textwidth]{svd/CONJ_TRUE.png}
%      \end{subfigure}
%         \caption{Convergence of Right Singular Vector to Left Eigenvector for powers of random $\A$ matrices of dimension 10. When the top two eigenvalues are conjugate pairs there is no convergence. When they are not, there is convergence.  }
%         \label{fig:svd_conj}
% \end{figure*}


\section{Detecting the dominant mode}
In this Section, we apply our formalism to time-series synthetically generated by a linear dynamical system of order $k$. 
% To develop intuition, we focus on small $k$ and assume that $\lambda_i\in  \R$.
%We postulate that the neuron attempts to infer a low-rank approximation to $\A$. A similar task is accomplished by a method called Dynamical Mode Decomposition (DMD) which finds an eigendecomposition of the matrix $\A$ learned from data. Unlike NMD, where eigenvalues are typically complex describing oscillatory dynamics, we focus on real eigenvalues describing transient rather than stationary processes. 

\subsection{Problem formulation} 
As stated in the introduction, we assume that we are in the vicinity of a hyperbolic fixed point of the dynamics of the lag vector. This means that the dynamics of the linearized system can be decomposed as the sum of $k$ exponentials with real exponents determined by the Jacobian of $f$ at the fixed point and coefficients determined by the initial conditions of the system. We also assume that an uncorrelated observation noise (see Fig.~\ref{fig:sample_series}):
\begin{equation}\label{eq:series_form}
    % x(t) = \sum_i^k x_i(t) +\eta(t)= \sum_i^k c_i e^{a_i t} +\eta(t).
    x(t) = \sum_i^k c_i e^{a_i t} +\eta(t).
\end{equation}
 %different exponents $a_i$ as well as the individual components $x_i(t)$. In terms of Fig.~\ref{fig:sample_series}, from the time-series given in the black curve, we would like to extract individual constituents. In particular, we are interested in 

% \begin{figure}[ht]
%     \centering
%     \includegraphics[width=0.5\textwidth]{finalfigures/sample_series.pdf}
%     \caption{ This figure should have three panels showing i) the three constituent exponents and noise ii) the sum of the three exponents with noise iii) extracted dominant exponent} %A sample time series and its $3$ constituent exponentials. In this example we have added a multiplicative Gaussian noise with $\sigma= 0.03$ to each constituent.}
%     \label{fig:sample_series}
% \end{figure}
\begin{figure*}[ht]
    \centering
    % \includegraphics[width=0.99\textwidth]{synth_figs/problem.png}
    \includegraphics[width=0.24\textwidth, trim=5 0 5 0]{finalfigures/traj4learn.pdf}
    \hfill
    \includegraphics[width=0.24\textwidth, trim=5 0 5 0]{finalfigures/learnedfilter.pdf}
    \hfill
    \includegraphics[width=0.24\textwidth, trim=5 0 5 0]{finalfigures/newseries.pdf}
    \hfill
    \includegraphics[width=0.24\textwidth, trim=5 0 5 0]{finalfigures/extractedseries.pdf}
    \caption{Problem formulation. (A) We generate a number of trajectories in the vicinity of a hyperbolic fixed point. Each trajectory is comprised of a number of growing and decaying exponentials and white noise (see panel C). (B) From these trajectories, we learn the top left eigenvector of the inferred $\A$ matrix. We use this as the time kernel of our proposed neuron. (C) We apply the time kernel to a previously unseen trajectory from the same dynamical system. (D) We extract the most dominant mode by convolving the time series with the computed time kernel. In this example, the presence of the dominant mode is clear from the projection starting at $t=1$. However, if we only look at the time series in (C), the presence of a growing mode would not be clear until around time $t=2$. Dynamics details: the time series is comprised of five different exponents with exponents given by $\{-1.5, -1, 0, 1, 1.5\}$. The time series is discretized with time-step 0.05.}
    \label{fig:sample_series}
\end{figure*}

In general, NMD extracts the exponents via an eigendecomposition, i.e. the eigenvalues of $\A$, $\lambda_i=e^{a_i}$. Since a neuron has a single output it can extract only one constituent. Because the greatest contribution to the future is given by the largest exponent (Figure~\ref{fig:sample_series}) we propose that a neuron learns the top left eigenvector of $\A$ and projects its lag vector, $\x_t$, onto this eigenvector which can then be used for prediction and control. 



% \subsection{Solution}
%The first key observation in this problem is that a time series of the form given in Eq.~\eqref{eq:series_form} is a \emph{linear} dynamical system of order $k$. The system is effectively described by its $k$ individual constituents, however, in the absence of this information, the same information can be derived from a lag-vector of length $n$:\footnote{Note that in general, the longer the lag-vector, the more linear a problem becomes. This is due to the fact that in a time series with fixed time-step, a Taylor expansion is a linear operation.} 
%\begin{equation}
 %   \x_{t+1} = \A \x_t\quad,\quad\x_t=[x(t-n+1) ... x(t)]^\top.
%\end{equation}

%The second key observation in this problem is that the constituent exponentials are eigenvectors of the time-translation operator $\A$. This is clear since, barring noise fluctuations, $x_i(t+1) = e^{a_i} x_i(t)$, hence the $i$'th constituent is an eigenvector with eigenvalue:
%\begin{equation}\label{eq:eigenvalues}
%    \lambda_i  = e^{a_i}.
%\end{equation}
% Because $\x_t$ and $\x_{t+1}$ in \eqref{dyn} are lag vectors $\A$ matrix has a so-called companion matrix form:
% \begin{align*}\label{eq:companion}
%     \A=
%     \begin{bmatrix}
%         0 & 1 &  &   \\
%          & 0 & 1 &   \\
%          &  & \ddots & \ddots &   \\
%          &  &  &  0 & 1 \\
%         c_1 & c_2 & \hdots & &c_k
%     \end{bmatrix}.
% \end{align*}
% A companion matrix can be diagonalized using the Vandermonde matrix and its inverse:
% \begin{align*}
%     \A=
%     \begin{bmatrix}
%         1&1&\cdots&1\\
%         \lambda_1&\lambda_2&\cdots & \lambda_k\\
%         \vdots& \vdots & & \vdots\\
%         \lambda_1^k & \lambda_2^k & \cdots & \lambda_k^k
%     \end{bmatrix}
%     \begin{bmatrix}
%         \lambda_1 & &\\ &\lambda_2 & \\ & &\ddots \\ & & &\lambda_k
%     \end{bmatrix}
%     \begin{bmatrix}
%         1&1&\cdots&1\\
%         \lambda_1&\lambda_2&\cdots & \lambda_k\\
%         \vdots& \vdots & & \vdots\\
%         \lambda_1^k & \lambda_2^k & \cdots & \lambda_k^k
%     \end{bmatrix}^{-1},
% \end{align*}
% where the eigenvalues, $\lambda_i$ are the roots of the characteristic polynomial with coefficients, $c_i$.  Then, the top left eigenvector, $v_1$, is given by the first row of the inverse of the Vandermonde matrix. 

% Whereas a scalar output encodes a projection of the input on only the top left eigenvector, its form depends on the other right eigenvectors of $\A$ because of their bi-orthogonality \eqref{DMD}. As the right eigenvectors corresponding to real eigenvalues are exponential, bi-orthogonality implies that the left eigenvector is multi-phasic, i.e. alternating sign.  %Eq.~\eqref{eq:eigenvalues}. The matrix on the left of the decomposition is called the Vandermonde matrix and as expected its columns are the right eigenvectors of $\A$. Therefore, in order to project onto the top eigenvector of the system, we need to compute the inverse of the Vandermonde matrix. In practice it is easier to simply look at the top left eigenvector of $\A$. 

\subsection{Dependence of the filter on the noise level}
In this section, we look at the effect of noise on the shape of the time kernel. We first discuss the effect of additive Gaussian noise analytically, we then verify this numerically on synthetically generated data. Previous work analyzing the application of NMD to data with observation noise have primarily focused on extracting the true eigenvalues and eigenvectors of the system, that is the dynamic modes of the dynamical system in the absence of noise~\cite{noisy_dmd,stable_noisy_dmd}. Here, we are interested in extracting the most dominant mode from noisy data, and as we will see (Fig.~\ref{fig:comp2david}), the eigenvector of the noiseless system is not a good candidate for this task.

\paragraph{Analytical calculation.} Let us assume that the time series given in Eq.~\eqref{eq:series_form} has additive white observation noise, that is we assume that $\X$ satisfies Eq.~\eqref{dyn} but we observe $\X_\e=\X + \eta$ with $\eta\sim \mathcal N(0,\epsilon^2)$. We denote the noiseless data and ground truth dynamics by $\X$ and $\A$, and call the data and the inferred dynamics in the presence of noise $\X_\e$ and $\A_\e$. When averaging over the different draws of the noise, we have $\E_\eta(\X_\e\X_\e^\top) = \X\X^\top + \epsilon^2 \I_k$ and $\E_\eta(\X_{\e+}\X_\e^\top) = {\X}_+\X^\top + \epsilon^2 \S$, where $\I_k$ is the identity matrix and $\S$ is the matrix of off-diagonal ones:
\begin{equation*}
    \S=
    \begin{bmatrix}
        0 & 1 &  &   \\
         & 0 & 1 &   \\
         &  & \ddots & \ddots &   \\
         &  &  &  0 & 1 \\
        0 & 0 & \hdots & &0
    \end{bmatrix}
\end{equation*}
Note that here $\E_\eta$ denotes averaging only over the noise draws and the averaging over the samples is performed in the product of data matrices ${\X}_+\X^\top$ and ${\X}_+\X^\top$. The putative $\A_\e$ matrix (that is the $\A$ matrix that minimizes the MSE objective with noisy data $\X_\e$) is now given by
\begin{multline}
    \label{eq:A-noisy}
    \A_\e = (\X_{\e+}\X_\e^\top)(\X_\e\X_\e^\top)^{-1}\\=(\A\X\X^\top+\epsilon^2\S)(\X\X^\top+\epsilon^2\I_k)^{-1}.
\end{multline}
That is, given the ground truth dynamics and noiseless data covariances (i.e. $\A$ and $\X\X^\top$), we can determine the inferred dynamics and eigenvectors in the presence of white noise. We provide further details of this computation in the supplementary materials Sec.~\ref{app:analytic}. While the eigenvalue equations are derived analytically, they are not solvable in closed form and we look at their numerical solutions for varying noise and exponents.
Here, we summarize the results and provide some intuition.

When looking at the eigenvalues of $\A_\e$ (Figure~\ref{fig:analytic figures0}), we see that as the noise level $\epsilon$ increases, the real parts slowly decrease and eventually cross zero. At the same time, the imaginary parts of the eigenvalues start at zero and grow when the real part becomes negative. This is intuitively understandable: when the increased amount of noise reduces the ability of the system to differentiate different modes, it starts replacing these with complex oscillatory modes.

Furthermore, at the noise value for which a mode's eigenvalue's real part crosses zero, we see that the top left eigenvector's number of phases decreases by one. This is due to the fact that the multi-phasic structure of the left eigenvectors are responsible for the bi-orthogonal structure of the left/right eigenvectors. 

\begin{figure*}[ht]
    \centering
    \includegraphics[height=4.95cm, trim=0 0 0 0, clip]{finalfigures/left_anal.pdf}
    \hfill
    \includegraphics[height=4.95cm, trim=0 0 0 0, clip]{finalfigures/eigs_real_anal.pdf}
    \hfill
    \includegraphics[height=4.95cm, trim=0 0 0 0, clip]{finalfigures/eigs_imag_anal.pdf}
    \hfill
    \includegraphics[height=4.9cm, trim=0 -26 0 2, ]{finalfigures/bar_anal.pdf}
    \caption{Analytically derived eigenvectors and eigenvalues for the lag $n=5$ for different noise levels. The $x$ axis for the eigenvalue plots denotes the eigenvalues sorted from largest to smallest real part. System with exponents $\{0.7, 0.2, -0.1, -0.4, -1.6\}$ and time-step equal to 1 for simplicity.}
    \label{fig:analytic figures0}
\end{figure*}

% \begin{figure}[h]
%     \centering
%     \includegraphics[height=3.6cm, trim=0 0 0 0, clip]{finalfigures/left_anal_v2.pdf}
%     \includegraphics[height=3.6cm, trim=0 0 0 0, ]{finalfigures/bar_anal.pdf}
%         \includegraphics[height=3.6cm, trim=0 0 0 0, clip]{finalfigures/eigs_complex_anal.pdf}
%     \includegraphics[height=3.6cm, trim=0 0 0 0, ]{finalfigures/complexbar_anal.pdf}
%     \caption{Top left eigenvector versus noise.}
%     % \label{fig:analytic figures}
% \end{figure}

% In the supplementary materials, we provide a Mathematica widget which demonstrates further the change of the eigenvalues and eigenvectors as one varies the exponents and the noise level.

\paragraph{Numerical verification.} Above we looked at the eigenvectors of $\A_\epsilon$, i.e. after averaging over different noise draws. Here, we look at these for specific noise draws, that is for specific simulations of linear dynamical systems. We looked at the response of the algorithm to synthetic data constructed from dynamics in the vicinity of a hyperbolic fixed point. We initiated a number of trajectories with different initial conditions and noise draws. Our algorithm finds the top eigenvector on the dynamics inferred from these trajectories. We then applied the learned eigenvector to a previously unseen trajectory. Figure ~\ref{fig:sample_series}(Right) shows the projection of this time series onto the largest left eigenvector and we see that we correctly recover the dominant exponential. For details of this experiment and further results and figures see supplementary materials Sec.~\ref{app:analytic}.



% \begin{figure}[h]
%     \centering
%     \includegraphics[width=\textwidth, trim=0 0 0 0, clip]{finalfigures/proj_noisy_comp.pdf}
%     \caption{This figure should have three columns - input signal, filter, and exponential output (both ground truth and acctual). Each row corresponds to different noise. Projected time series. The dashed red line denotes the ground truth for the growing exponential mode.}
%     \label{fig:synth_performance}
% \end{figure}
% \begin{figure}[h]
%     \centering
%     \includegraphics[width=0.32\textwidth]{fig2/series.png}
%     \includegraphics[width=0.32\textwidth]{fig2/filter.png}
%     \includegraphics[width=0.32\textwidth]{fig2/proj.png}
%     \caption{Here we see the shape of the filter which extracts the fast growing exponential out of the time series. We plot the mean and standard deviation over 10,000 noise draws.}
%     \label{fig:analytic figures}
% \end{figure}

Figure~\ref{fig:synth_results} shows the effect of various parameters such as the noise, the system order, and the lag vector length on the shape of the filter. We see that adding noise decreases the number of the phases as expected (Fig.~\ref{fig:synth_results}A). Similarly as we decrease the system order $k$, the number of phases decrease  (Fig.~\ref{fig:synth_results}B) as expected from the bi-orthogonality argument given at the end of Sec.~\ref{sec:DMD}.  Figure~\ref{fig:synth_results}C shows that with a finite amount of noise, as the length of the lag vector is increased, the number of phases can slowly increase. This is due to the fact that with longer lag vectors, there is higher noise tolerance. Indeed as we take the noise values to zero, we see that the lag vector length no longer affects the number of phases in the filter. Increasing the length of the lag vector merely stretches the filter shape (Fig.~\ref{fig:vslength_lownoise} in the supplementary materials).

\begin{figure*}[ht]
     \centering
     \begin{subfigure}[b]{0.32\textwidth}
         \centering
        %  \fbox{
         \includegraphics[width=\textwidth, trim=0 1 13 10, clip]{finalfigures/vsnoisev2.pdf}
        %  }
         %\caption{Filter shape vs. noise}
         \label{fig:synth_vs_noise}
     \end{subfigure}
     \hfill
     \begin{subfigure}[b]{0.32\textwidth}
         \centering
         \includegraphics[width=\textwidth, trim=0 1 13 10, clip]{finalfigures/vsorder.pdf}
         %\caption{Filter shape vs. system order}
         \label{fig:synth_vs_order}
     \end{subfigure}
     \hfill
     \begin{subfigure}[b]{0.32\textwidth}
         \centering
         \includegraphics[width=\textwidth, trim=0 1 13 10, clip]{finalfigures/vslength.pdf}
         %\caption{Filter shape vs. lag vector length}
         \label{fig:synth_vs_length}
     \end{subfigure}
        \caption{Dependence of the filter shape for the rank 5 system in Fig.~\ref{fig:sample_series} on system parameters: (A) noise amplitude, (B) system order (the number of exponential constituents, $k$), and (C) length of the lag vector, $n$. In (B) and (C) the noise is fixed at $0.01\%$.  In (B), to achieve a comparable system of lower rank $k$, we keep the top $k$ exponents of the system.}
        \label{fig:synth_results}
\end{figure*}

% \section{Concatenated synthetic data}
% Here, we consider synthetic datasets that bridge the simplified example of the previous section with a natural stimulus that will be considered in the next. Specifically, we consider the scenario comprised of concatenation of different segments, where each segment is a time series given by the same overall linear dynamical system described by $\A$, but with different initial conditions (Figure~\ref{fig:concat_full} A). This is an approximation of a linear dynamical system which is occasionally perturbed by impulses in a randomly sampled direction.
% % (similar to a  ball that is occasionally kicked at a random angle or electron in EM field? electron in EM field can have exponential motion but ball does not).

% In this case, we find that naively applying our method from the previous section does not work well: the top eigenvalue of the $\A$ matrix when fit on the entire series is a poor estimate of the true eigenvalue of the dynamical system (Figure ~\ref{fig:concat_full}C). Because the stitching of different exponential regions over a whole dataset is best fit by a constant, the top exponent is close to 1, Figure ~\ref{fig:concat_full}C. To circumvent this problem, one approach would be to slice the data into different linear dynamics regions and fit them separately. However, this is not possible if the location of transition regions is not known a priori.

% Here, we take a different approach, inspired by neuronal physiology. We perform the NMD over a sliding window across the entire time series, and then average the resulting top left-eigenvector from each window to form our final estimate of the filter. Figure ~\ref{fig:concat_full}B shows the top eigenvalue plotted in time as the window slides over the series. Most of the time we have a strong estimate of the top eigenvalue. It is only when the window passes through the transition regions that there is significant error in our estimate. The eigenvalues produced in these transition regions average out leaving us with a strong estimate of the top eigenvalue, Figure ~\ref{fig:concat_full}C. We will see in the next section that this approach also leads to a competitive fit to the experimental data. 



% \begin{figure*}[ht]
%      \centering
%      \hfill
%     %  \begin{subfigure}[b]{0.32\textwidth}
%         %  \centering
%          \includegraphics[height=4.3cm]{synth_figs/concat.png}
%         %  \label{fig:concat_series}
%     %  \end{subfigure}
%      \hfill
%     %  \centering
%     %  \begin{subfigure}[b]{0.32\textwidth}
%         %  \centering 
%          \includegraphics[height=4.3cm]{synth_figs/eigval_full.png}
%         %  \label{fig:concat_full}
%     %  \end{subfigure}
%      \hfill
%     %  \centering
%     %  \begin{subfigure}[b]{0.32\textwidth}
%         %  \centering
%          \includegraphics[height=4.3 cm]{synth_figs/esti_ev.png}
%       \hfill
%     %  \end{subfigure}
%         \caption{Performance of the eigen-decomposition approach on the full concatenated exponentials dataset and of the sliding-window version. (A) A typical concatenated exponentials time series. Time windows overlapping transition regions are shaded. (B) Top eigenvalue obtained by the sliding-window approach discover the growing exponential in the linear segments but fail in the transition regions (shaded). (C) While the top eigenvalue of $\A$ computed on the full dataset does not discover the growing exponentials, averaging recovers the top eigenvalue from growing exponentials without a priori information about the location of transition regions.}
%         \label{fig:concat_full}
% \end{figure*}



% \subsection{Concatenated logistic functions}

% Next, we consider a series of transitions described by sums of logistic functions (ADD FIGURE). We can think of the sums of logistics as a idealized approximation for transitions between different light intensities. In this case, the time series is given by
% \begin{align}
%     x(t)=\sum_i\frac{c_i}{1+e^{-\lambda_i(t-s_i)}}+\eta(t).
% \end{align}
% For a given $i$ and $t\ll s_i$,
% \begin{align*}
%     \frac{c_i}{1+e^{-\lambda_i(t-s_i)}}\approx c_ie^{\lambda_i(t-s_i)},
% \end{align*}
% and for $t\gg s_i$,
% \begin{align*}
%     \frac{c_i}{1+e^{-\lambda_i(t-s_i)}}\approx c_i(1-e^{-\lambda_i(t-s_i)}).
% \end{align*}
% In particular, we can locally approximate $x(t)$ as a sum of exponentials and apply NMD to extract relevant modes.

\subsection{Comparison with optimal projection under noisy observations}

To verify that in the presence of noise our algorithm provides a sensible approximation of the dominant dynamical mode, in this section we discuss a noise-optimal projection alternative and provide empirical comparisons.

Let $\x_t$ be a time series and suppose our goal is to extract the dominant exponential mode by projecting $\x_t$ onto the filter given by $\v$ (derived as the top eigenvector of \eqref{eq:gen_eig_prob} when there is no observation noise). We refer to $\v$ as the \emph{noiseless filter}. Suppose, however, we do not have access to $\x_t$, but only a noisy observation $\y_t=\x_t+\n_t$, where $\n_t$ is isotropic Gaussian noise with variance $\alpha^2$. What is the optimal projection of $\y_t$? Consider the optimization problem
\begin{multline*}
    \min_{\w} E[|\w^\top\y_t-\v^\top\x_t|^2]=\\
    % =\min_\w E[|(\w-\v)^\top\x_t+\w^\top\n_t|^2]\\
    % =\min_\w E[|(\w-\v)^\top\x_t|^2]+\alpha^2\w^\top\w\\
    \min_\w (\w-\v)^\top\C_{xx}(\w-\v)+\alpha^2\w^\top\w,
\end{multline*}
where we have used the fact that $\n_t$ is mean zero isotropic noise independent of $\x_t$. Solving for the optimal $\w$ we find:
\begin{align}
    \w=(\C_{xx}+\alpha^2\I)^{-1}\C_{xx}\v.
    \label{eq:noise_optimal}
\end{align}
% Note that this is strongly dependent on $\C_{xx}$, which is difficult to interpret in our setting because $\x_t$ is not stationary. If we assume that $\C_{xx}$ is identity, then increasing noise (i.e., $\alpha$) simply scales the optimal project, but does not change the optimal \textit{direction}.
In order to apply this method, we need to know the noise variance $\alpha$ and also the noise-free projector $\v$ and covariance $\C_{xx}$. Computing these noise-free quantities is challenging, especially in a problem where the dynamics might be changing over time. Furthermore, computing the full covariance matrix $\C_{xx}$ requires more samples than just computing the top subspace in our generalized eigenvalue problem (Eq.~\eqref{eq:gen_eig_prob}). In practice, we find that even if we know $\alpha$ and $\v$, our method gives comparable performance to this noise-optimal solution (Fig.~\ref{fig:comp2david}). Furthermore, our method provides a better estimate of the most dominant mode than the non-adaptive method which uses the true (i.e. noiseless) eigenvector. In practice, in order to compute the noiseless filter $\v$, we set the noise variance $\alpha=10^{-6}$.

\begin{figure}
    \centering
    \includegraphics[width=0.45\textwidth]{finalfigures/comp2david.pdf}
    \caption{Comparison of the performance with respect to the mean square error between the recovered dominant exponential and the ground truth. The `noise optimal' solution is given by the filter in Eq.~\eqref{eq:noise_optimal} and `noiseless filter' refers to using the filter compted in the absence of noise but applied to noisy observations. The system is the same as in Fig.~\ref{fig:sample_series}. }
    \label{fig:comp2david}
\end{figure}

% \section{Natural input Data}

% As an application of our formalism, we apply our algorithm to  in vivo electrophysiological recordings in the early visual system of {\it Drosophila} \citep{zheng2006feedback,zheng2009network}. Photoreceptors transduce the incoming light and provide synaptic input to the large monopolar cells (LMC). Unlike the vertebrate retina, LMCs do not integrate over space and receive essentially scalar input from photoreceptors. Both photoreceptors and LMCs are non-spiking graded potential neurons~\footnote{This dataset is inherently non-linear and does not strictly fall under the purview of our assumptions. However, as mentioned in the introduction, NMD methods have been applied to non-linear systems~\cite{schmid2010dynamic,kutz2017dynamic}.}. 

% \begin{figure*}[ht]
%      \centering
%      \begin{subfigure}[b]{0.32\textwidth}
%          \centering
%          \includegraphics[width=\textwidth]{realdata_figs/stim.png}
%      \end{subfigure}
%      \hfill
%      \centering
%      \begin{subfigure}[b]{0.32\textwidth}
%          \centering
%          \includegraphics[width=\textwidth]{realdata_figs/ph.png}
%      \end{subfigure}
%      \hfill
%      \begin{subfigure}[b]{0.32\textwidth}
%          \centering
%          \includegraphics[width=\textwidth]{realdata_figs/lmc.png}
%      \end{subfigure}
%         \caption{(A) Light stimulus, (B) Photoreceptor response and LMC input, and (C) LMC neuron responses. The shaded regions represents one standard deviation computed over 100 repetitions of the  light stimulus.}
%         \label{fig:LMC-input-output}
% \end{figure*}

% Our dataset \cite{zheng2006feedback,zheng2009network} was obtained using in-vivo whole cell recordings from photoreceptors and LMCs stimulated with naturalistic time-varying light stimulus \cite{van1997processing} (Figure~\ref{fig:LMC-input-output}). Here we focus on the transformation between two time series: the output of the photoreceptors, i.e. input to the LMC, $x(t)$, and the corresponding output of the LMC neurons, $y(t)$. Our goal is to predict the output, $y(t)$, of LMC in response to the input, $x(t)$. 

% The standard approach to this problem is using supervised learning of the optimal linear filter  $\k=[k_1,...,k_n]^\top$ that minimizes the error as in: 
% \begin{align*}
% \min_\k\|\y-\k^\top\X\|^2
% \end{align*}
% where $\y=[y(n),...,y(T)]$. We solve this least squares problem using pseudoinverse as is common for linear regression by using the data from the training set and measure the error on the test set. By contrast our method proceeds in a unsupervised way,

% %The experimental task is to model the algorithmic activity of the LMC neurons by predicting the LMC response, $y(t)$, from the photoreceptor response, $x(t)$ where $t \in [0,T]$. We formalize this objective by splitting the data into train and test sets. The goal is to use the training set to learn a time kernel $k(t)$ of dimension $d$ such that we minimize the error as in: $\|k(t) * x(t) - y(t)\|^2$.

% %The classical way of approaching this task is a supervised method which fits a time kernel using a linear least squares method. To do this we begin by building a $d \times T$ lag matrix, $\X$, whose columns $d$ dimensional lag vectors of the time series $x(t)$:
% %\begin{align*}
% %    \mathbf{X}=
% %    \begin{bmatrix}
% %        x_d      & x_{d+1}  & \hdots & x_{T-d}  \\
% %        x_{d+1}  & x_{d+2} & \hdots & \vdots   \\
% %        x_{d+2}  & x_{d+3} & \hdots & x_{T-2}  \\
% %        \vdots  & \vdots   & \hdots & x_{T-1}  \\
% %        x_{2d}   & x_{2d+1}   & \hdots & x_{T}
% %    \end{bmatrix}.
% %\end{align*}

% %By this construction the vector matrix product $\k\X$ represents a discrete time convolution of $k(t)$ with $x(t)$. Thus by setting $\k\X=\y$ we can we can solve for $\k$ using a least squares method. The solution is the optimal linear time kernel on the training data.

% %By contrast our method proceeds in a unsupervised way, learning a time kernel from only $x(t)$. We begin by implicitly learning an $\A$ matrix which gives the linear dynamics of $x(t)$ in the sense that $\x_{t+1}=\A \x_t$. We then find the left eigenvector with the largest eigenvalue which we use as our time kernel. We accomplish this by first constructing two $d \times T-1$ lag matrices, where $\X_0$ is built from $x_0 \dots x_{T-1}$ and $\X_{+}$ is built from $x_1 \dots x_{T}$. We then solve the following generalized eigenvalue problem for left eigenvectors and eigenvalues $\w$ and $\lambda$ with a regularization parameter $\beta$:
% %$$
% %\lambda \w^\top  \left(\X_{0} \X_{0}^\top + \beta \I \right) = \w^\top  \left(\X_{+} \X_{0}^\top \right)
% %$$
% %Adding regularization in this way is equivalent to solving for $\A$ via ridge regression on $\X_{0}$ and $\X_{+}$, and then finding the left eigenvectors of $\A$. That is, if $\w$ and $\lambda$ are left eigenpairs of $\A$ then solving for $\A$ via ridge regression gives,
% %\begin{align*}
% %A = \left(\X_{+} \X_{0}^\top\right) \left(\X_{0} \X_{0}^\top + \beta \I \right)^{-1} \\
% %\w^\top A = \w^\top \left(\X_{+} \X_{0}^\top\right) \left(\X_{0} \X_{0}^\top + \beta \I \right)^{-1} \\
% %\lambda \w^\top = \w^\top \left(\X_{+} \X_{0}^\top\right) \left(\X_{0} \X_{0}^\top + \beta \I \right)^{-1} \\
% %\lambda \w^\top \left(\X_{0} \X_{0}^\top + \beta \I \right) = \w^\top \left(\X_{+} \X_{0}^\top\right)
% %\end{align*}

% %\subsection{LMC Results}


% We apply the method described in the previous section to our real world dataset with one modification: instead of computing the top eigenvector of the entire data matrix, we look at the eigen-decomposition of a sliding window of a specific length and then average the resulting top eigenvector for the different windows. We do this modification in order to increase the biological plausibility of our algorithm Similarly to the supervised method, we split the inputs and outputs into train and test sets in order to build a model and then measure its generalization ability, Figure~\ref{fig:train-test-pred}. Our plots in Figure~\ref{fig:LMC-train-test} show the results of these experiments as we vary the length of the lag vector. Additionally, we perform 100 different runs, sampling the test set from a different portion of the original time series. This allows us to plot the mean and one standard deviation for each value of the lag. We see from these runs that our method easily out performs the identity baseline. 
% \begin{figure*}[ht]
%     \includegraphics[width=.42\textwidth]{finalfigures/learnfilter.pdf}
%     \hspace{10pt}
%     \includegraphics[width=.42\textwidth]{finalfigures/train_test_pred.pdf}
%     \caption{Results of our method on neuronal data. (A) Top left eigenvector learned on LMC input. (B) LMC output predicted by the convolution of the top left eigenvector with train and test segments of input data.}
%     \label{fig:train-test-pred}
% \end{figure*}

% On the training set our unsupervised method performs worse than the supervised method. This would be expected because the former produces a single average filter while the latter gives the optimal linear filter for the data on which it is trained. What is remarkable about our method is its performance on the test set. We perform on par with the optimal linear method for lags of $30$ and $40$ms. Interestingly, this timescale is within experimental estimates of neuronal membrane time constants~\cite{10.5555/1137840}. Moreover the qualitative shape of the filter which we extract (Fig.~\ref{fig:train-test-pred}A) matches that from experimental observations as well as the synthetic experiments discussed previously. This provides strong evidence of our method as a plausible model of a neuron's computational behavior. Interestingly, we find that performing a windowed computation with a window length given by biological time scales provides a better fit for these empirical observations.

% \begin{figure*}[ht]
%          \includegraphics[width=.42\textwidth]{finalfigures/train.pdf}
%          \hspace{10pt}
%          \includegraphics[width=.42\textwidth]{finalfigures/test.pdf}
%         \caption{LMC output prediction error as a function of lag vector length for three methods: Identity transformation (blue), supervised optimal linear filter computed by  linear regression (orange), our sliding window eigen-projection method (green). The error bars represent one standard deviation, where the variance is incurred by varying the test set as described in the supplement. (A) Training error. (B) Test error.}
%         \label{fig:LMC-train-test}
% \end{figure*}

% % \begin{figure}[h]
% %      \centering
% %      \begin{subfigure}[b]{0.32\textwidth}
% %          \centering
% %          \includegraphics[width=\textwidth]{lmc_finalfigures/BG4_eigfilter.png}
% %          \caption{Filter shape for eigen-projection for optimal choice of $\beta$}
% %      \end{subfigure}
% %      \hfill
% %      \begin{subfigure}[b]{0.32\textwidth}
% %          \centering
% %          \includegraphics[width=\textwidth]{lmc_finalfigures/BG4_sub_linfilter.png}
% %          \caption{Filter shape from linear for slightly sub-optimal choice of $\beta$}
% %      \end{subfigure}
% %      \hfill
% %      \begin{subfigure}[b]{0.32\textwidth}
% %          \centering
% %          \includegraphics[width=\textwidth]{lmc_finalfigures/BG4_linfilter.png}
% %          \caption{Filter shape from linear for optimal choice of $\beta$}
% %      \end{subfigure}
% %         \caption{The shape for the time kernels for BG4}
% %         \label{fig:LMC-filters}
% % \end{figure}
% % \begin{figure}[h]
% %      \centering
% %      \includegraphics[width=0.5\textwidth]{lmc_finalfigures/all_eigfilters.png}
% %      \caption{Filter shapes for eigen-projection for optimal choice of $\beta$ across all noise levels}
% %      \label{fig:LMC-filters-noise}
% % \end{figure}

% % \subsection{Neural data (Shiva/Tibi?)}

% \vspace{10pt}

\section{Discussion}
We propose to model a neuron on an algorithmic level as performing a NMD on its input. Specifically, the temporal filter of a neuron is given by the top left eigenvector of the generalized eigenproblem formulated in terms of covariances of lag vectors. The neuron outputs the fastest growing mode present in the input thus predicting a future trend.

How could a biological neuron find the top eigenvector of the generative matrix? One possibility is for a neuron to solve the generalized eigenproblem \eqref{eq:gen_eig_prob} using a power iteration. Specifically, this would require two types of history-dependent active conductances (ion channels) which encode input covariances, $\mathbf{X_+X^\top}$ and $\mathbf{XX^\top}$, with opposite signs (depolarizing and hyperpolarizing). Then the solution of \eqref{eq:gen_eig_prob} would be given by the voltage that balances the two currents. The output of a neuron would reflect such voltage, thus projecting the input on the top left eigenvector of the transition matrix.

Whereas a single neuron computing a scalar signal can represent only one eigenmode, multiple neurons may extract multiple modes (not just the top eigenmode). For each neuron to represent a different mode their output must be decorrelated, for example, by the lateral inhibitory connections. The strengths of such connections can be adjusted using biologically plausible local learning rules as has been shown by some of the authors previously in the context of extracting eigenmodes of the multichannel covariance \cite{pehlevan2015hebbian,pehlevan2017similarity,pehlevan2019neuroscience}. Also, we proposed a framework for deriving multichannel neural networks for solving symmetric generalized eigenvalue problems \cite{PRXLife.1.013008} and the approach can potentially be extended to solve non-symmetric generalized eigenvalue problems. We anticipate that a similar approach can be applied to extracting eigenmodes of the time-lagged dynamics using local learning rules.

Because the matrix $\A$ in Eq.\eqref{dyn} is generically non-normal, its eigenvectors are not guaranteed to be orthogonal. In practice, this means that they are not robust to perturbations to the elements of the matrix $\A$. This requires particular care to ensure that the eigenvalues are sufficiently distinct as determined by the pseudospectrum of $\A$ \cite{TrefethenEmbree+2020}.  

Neuronal projection of its input onto the subspace corresponding to the fastest growing mode has an interpretation in terms of the phase portrait of the generative dynamical system. Close to a hyperbolic fixed point the dynamics are linear and the fastest growing mode would correspond to the unstable invariant subspace which evolves into an attracting manifold away from the fixed point. Therefore, the sign of the projection onto the unstable subspace predicts along which unstable manifold the future trajectory of the system will develop beyond the linear regime. If the neuronal output is rectified, it effectively assigns the trajectory to a particular future state. Output of a layer of such rectifying neurons becomes a latent vector variable that can in turn be an input to the next layer. This opens a path to stacking the layers of such neurons to discover more and more abstract latent variables. 

% \section*{Broader impact}

% We do not foresee any potential societal harm caused by our work. We hope our contribution moves forward the understanding of neural computations.

\bibliography{references.bib}
\bibliographystyle{unsrt}
% \clearpage

% \section{Bibliography}

% Zheng L, de Polavieja GG, Wolfram V, Asyali MH, Hardie RC \& Juusola M (2006) Feedback network controls photoreceptor output at the layer of first visual synapses in Drosophila. J Gen Physiol, 127(5), 495-510.
% Zheng L, Nikolaev A, Wardill TJ, O'Kane CJ, de Polavieja GG \& Juusola M (2009) Network adaptation improves temporal representation of naturalistic stimuli in Drosophila eye: I dynamics. PLoS ONE, 4(1).


%%%%%%%%%%%%%%%%%%%%%%%%%%%%%%%%%%%%%%%%%%%%%%%%%%%%%%%%%%%%

\clearpage
\appendix

\begin{center}\LARGE{\textbf{Supplementary Materials}}\end{center}
\vspace{3pt}
This is the supplementary materials section for the Physical Review Research submission titled ``Neuronal Temporal Filters as
Normal Mode Extractors''.

% \section{Numerical Experiments — Synthetic Data}\label{app:synth}

% As discussed in the paper we  produce a synthetic time series $x$ as
% \begin{equation}\label{eq:synth_gen_model}
%     x(t) = \sum_i^k c_i e^{a_i t} +\eta(t).
% \end{equation}

% Figure \ref{fig:sample_series} is produced using the following values of $a=[-0.8, -0.15, 1.5]$ and $c=[1.0, 0.2, 0.01]$ over a time interval $T=[0,3]$ with at time step of $d_t=0.1$. These values are selected to be demonstrative of a series with a clear decaying and growing exponential components over the time interval. Other figures use different sets of parameters. For the details of parameters for each experience see the included code. 

% In the case of concatenated synthetic data, we take $N$ many exponential time series as discussed above. These are then concatenated horizontally in time, with each series shifted upward by the previous series' endpoint. Figure 5 is generated from 8 time series with the same $c$ and $a$ values given above where each is then multiplied by a random sign, -1 or 1.

% The concatenated data motivates the use of a windowed NMD method. In this case, the window size is taken to be three times the number of components in the lag vector. Given a time series $x$ with $T$ many points and some window size $w$, we take a contiguous slice of the time series from $x_{t-w}$ to $x_{t}$. We then solve for an $\A$ matrix as given in Eq. \ref{msec} and extract the top eigenvector $v_t$. This process is repeated for each time point in $[w,T]$ resulting in a collection of $T-w$ eigenvectors over which the mean eigenvector $v_m$ is computed. The vector $v_m$ is then taken to be the approximate filter.

% On synthetic data using a windowed method it is possible that the covariance matrix $\X\X^\top$ used to compute $\A$ ends up not being numerically invertible. In order to avoid this issue add a small perturbation, $\beta$, along the diagonal of $\X\X^\top$ where $\beta= 10^{-5}$ such that we have,
% \begin{equation}
%  \a = (\x_+\X^\top)(\X\X^\top + \I\beta)^{-1}
% \end{equation}
% Where $\a$ is the bottom row of the companion matrix $\A$.

\section{Details of analytic arguments}\label{app:analytic}

Starting with the equation~\eqref{eq:A-noisy} describing the effect of noise:
\begin{equation}
    \A = (\X_+\X^\top)(\X\X^\top)^{-1}=(\A_0\X_0\X_0^\top+\epsilon^2\S)(\X_0\X_0^\top+\epsilon^2\I_k)^{-1}.
\end{equation}
we can derive the noisy version of the dynamics if we know the covariance structure of the non-noisy data $\X_0\X_0^\top$. It is intuitively understandable why this structure is necessary since the effect of noise depends not just on the eigenvalues and eigenvectors of $\A$, it also depends on the magnitudes of each individual mode. Specifically, if we assume the generative model in Eq.~\ref{eq:series_form}, it is not just the $a_i$ that determine the noisy dynamics, the $c_i$ are also important and this information is not present in $\A$ but is available in  $\X_0\X_0^\top$.

For plots given in Fig.~\ref{fig:analytic figures0}, we computed the $\X_0\X_0^\top$ using time series of the form in Eq.~\ref{eq:series_form}. We take the length of the time series, to be just long enough to compute a full rank covariance matrix.

% These arguments are implemented in a Mathematica widget included in the supplementary material files. In this widget the $a_i$ and $c_i$ are as in Eq.~\eqref{eq:series_form} and $\epsilon$ controls the amount of noise in log scale, that is the variance of the noise is equal to $10^\epsilon$.


\section{Other Figures}\label{app:figs}
Figure~\ref{fig:vslength_lownoise} shows the dependence of the filter shape as we vary the lag vector length for a small amount of noise. We see that increasing the lag vector length no longer affects the number of phases (cf. Fig.~\ref{fig:synth_results}C).

\begin{figure}[t]
    \centering
    \includegraphics[width=0.3\textwidth]{finalfigures/vslength_lownoise.pdf}
    \caption{The dependence of the filter shape vs the length of the lag vector at low noise values ($0.00001\%$).}
    \label{fig:vslength_lownoise}
\end{figure}


% \begin{figure*}
%     \centering
%     \includegraphics[width=\textwidth]{finalfigures/widget.png}
%     \caption{Demonstration of the Mathematica widget included in the supplementary material files which demonstrates the behavior of the eigenvalues and eigenvectors as a function of noise.}
%     \label{fig:widget}
% \end{figure*}

% \section{Numerical Experiments — Real Data}\label{app:real}
% The photoreceptor and LMC data are sampled at $10$kHz over a time period of 1 second, resulting in a time series of $10000$ points. This dataset is generously provided by \citep{zheng2006feedback,zheng2009network}. Before any analysis is done, the data is downsampled by a factor of 20 resulting in a time series of $500$ points with a time step of $2.0$ milliseconds. There are $100$ recordings generated from repetitions of the same light stimulus. The data we work with is taken as the mean of these repetitions.

% The train-test split is made in an 80/20 ratio resulting in a time series of 400 and 100 points for train and test respectively. Different train/test splits are made by shifting the starting point of the test set 100 times uniformly spaced throughout the time series. Given some starting point $i$, the test set is constructed from contiguous time points $x_i, i \in [i,i+100)$. The train set is then constructed from the non-contiguous time points $x_i, i \in [0,i) \cup [i+100,500)$. This can be seen in Fig \ref{fig:ttmid}. Clearly, this can result in a discontinuity in the training data. But any issues can be avoided if the Hankel matrix is constructed such that no lag vector crosses the discontinuity between $x_i, i \in (i,i+100]$.
% It is the shifting of the test set that is the source of variance which generates the error bars in Fig \ref{fig:LMC-train-test}.

% \begin{figure}[b]
%     \centering
%     \includegraphics[height=5cm]{finalfigures/train_test_mid_split.pdf}
%     \caption{Here we see an example of how the X data (photoreceptor response) is split when the test set is in the middle. The train set Hankel matrix is then constructed such that no lag vector cross the split in the train set.}
%     \label{fig:ttmid}
% \end{figure}

% The photoreceptor and LMC data are not normalized. Thus some methods may output their prediction on the wrong scale or bias. To resolve this issue, we allow each method to learn a scale parameter, $a$, and shift parameter $b$ during the training phase. That is given the target single $y$ and the predicted signal $\hat{y}$ we solve:
% \begin{equation}
% \min_{a,b} |y - a\hat{y} + b|^2
% \end{equation}
% The $a$ and $b$ parameters computed during the training phase are then used to scale and shift a method's output during the test phase.


% \section{Code and compute requirements}\label{app:code}
% We have attached a zip file containing python notebooks as well as the data \citep{zheng2006feedback,zheng2009network} used for creating the figures in this paper. Additionally, we include the Mathematica notebook containing the widget described above.

% In terms of compute requirements, all experiments were run on a 2016 Macbook Pro with a 2.9 GHz Quad-Core Intel Core i7 and 16GBs of RAM. No experiment took longer than 4 minutes to run start to finish. We do not anticipate compute being an issue in running the code.

\section{Inverse of a Vandermonde matrix}\label{app:vandermonde}

Recall that in the noiseless case (i.e., $\epsilon=0$), the matrix $\A$ has eigendecomposition $\A=\V\Lam\V^{-1}$, where $\V$ is a Vandermonde matrix. In this section we show that the top left eigenvector of a Vandermonde matrix of rank $n$ changes sign $n-1$ times. Consider the Vandermonde matrix
\begin{align*}
    \V:=\begin{bmatrix}
        1 & 1 & \cdots & 1 \\
        \lambda_1 & \lambda_2 & \cdots &\lambda_n\\
        \vdots & \vdots & & \vdots\\
        \lambda_1^{n-1}&\lambda_2^{n-1}&\cdots&\lambda_n^{n-1}
    \end{bmatrix}
\end{align*}
where $\lambda_1>\cdots>\lambda_n>0$. Rawashdeh \citep{rawashdeh2019simple} showed that the inverse of the Vandermonde matrix $V$ is given by
\begin{align*}
    (\V^{-1})_{ij}=(-1)^{i+j}S_{n-j,i}\frac{\prod_{k<l\text{ s.t.\ }l,k\ne j}^n(\lambda_k-\lambda_l)}{\prod_{k<l}^n(\lambda_k-\lambda_l)},
\end{align*}
where
\begin{align*}
    S_{j,k}:=\sum_{1\le i_1<\cdots<i_k\le n,i_\ell\ne j}\prod_{m=1}^k\lambda_{i_m}>0.
\end{align*}
Importantly, the sign of the $(1,j)$\textsuperscript{th} entry is positive (resp.\ negative) if $j$ is odd (resp.\ even). Therefore, the top left eigenvector changes sign $n-1$ times.


\section{Imaginary Parts in Eigenvalues}
In the main text we only examine a time series composed of real exponents. Here we examine the case where exponents have imaginary parts which creates oscillatory behavior in the training and test time series. We reproduce Figure \ref{fig:sample_series} where we infer the $\A$ matrix from training trajectories, extract the top left eigenvector and use it to extract the dominant mode from an unseen time series. In Figure \ref{fig:imag_nolarge_same} (Row A) we add an imaginary part in all trajectories except the one with the largest eigenvalue. We see that we are still able to accurately extract the largest mode of the time series. In Figure \ref{fig:imag_nolarge_same} (Row B) we add an imaginary part in all trajectories including the one with the largest eigenvalue. Here we fail to accurately reconstruct the largest mode, although we get the qualitative behavior correct. In Figure \ref{fig:imag_nolarge_same} (Row C) we repeat the test in Row B but in addition we lower the noise level. Under these circumstances we see that we are able to accurately reconstruct the largest mode, even when it contains an imaginary part. The specific parameters are given in Table \ref{tbl:imgtable}.


\begin{table}
  \centering
  \begin{tabular}{l|l|l}
    \toprule
    \cmidrule(r){1-2}
    Row  & Exponent Imaginary Part  & noise \\
    \midrule
    A & $\{-5+0.2i, -1.5+7i, 1i, 1.5+0.1i, 5.001\}$ &  $0.05$ \\
    B & $\{-5+0.4i, -1.5+0.2i, 0.5i, 1.5+0.5i, 5.001+6i\}$ &  $0.05$ \\
    C & $\{-5+0.4i, -1.5+0.2i, 0.5i, 1.5+0.5i, 5.001+6i\}$ &  $0.0001$ \\
    \bottomrule
  \end{tabular}
  \caption{Experimental parameters for Figure \ref{fig:imag_nolarge_same}}
\label{tbl:imgtable}
\end{table}

\begin{figure*}[t]
    \centering
    \includegraphics[width=0.24\textwidth, trim=5 0 5 0]{imaginary_figures/A1.pdf}
    \hfill
    \includegraphics[width=0.23\textwidth, trim=5 0 5 0]{imaginary_figures/B1.pdf}
    \hfill
    \includegraphics[width=0.24\textwidth, trim=5 0 5 0]{imaginary_figures/C1.pdf}
    \hfill
    \includegraphics[width=0.24\textwidth, trim=5 0 5 0]{imaginary_figures/D1.pdf}
    \\
    \includegraphics[width=0.24\textwidth, trim=5 0 5 0]{imaginary_figures/A3.pdf}
    \hfill
    \includegraphics[width=0.23\textwidth, trim=5 0 5 0]{imaginary_figures/B3.pdf}
    \hfill
    \includegraphics[width=0.24\textwidth, trim=5 0 5 0]{imaginary_figures/C3.pdf}
    \hfill
    \includegraphics[width=0.24\textwidth, trim=5 0 5 0]{imaginary_figures/D3.pdf}
    \\
    \includegraphics[width=0.24\textwidth, trim=5 0 5 0]{imaginary_figures/A2.pdf}
    \hfill
    \includegraphics[width=0.23\textwidth, trim=5 0 5 0]{imaginary_figures/B2.pdf}
    \hfill
    \includegraphics[width=0.24\textwidth, trim=5 0 5 0]{imaginary_figures/C2.pdf}
    \hfill
    \includegraphics[width=0.24\textwidth, trim=5 0 5 0]{imaginary_figures/D2.pdf}

        \caption{We reproduce Figure \ref{fig:sample_series} with imaginary parts in the exponents. (A) Imaginary parts in all trajectories except the one with the largest eigenvalue. (B) Imaginary parts in all trajectories including the one with the largest eigenvalue. (C) Imaginary parts in all trajectories including the one with the largest eigenvalue, but lower noise level.}
            \label{fig:imag_nolarge_same}
\end{figure*}



\clearpage



\end{document}